\documentclass[11pt,a4paper]{article}
\usepackage{fullpage}
\usepackage[utf8]{inputenc}
\usepackage[T1]{fontenc}
\usepackage[french]{babel} \frenchbsetup{og=«,fg=»}
\usepackage{graphicx}
\usepackage{amsthm}
\usepackage{hyperref}
\usepackage{amsmath}
\usepackage{amssymb}
\usepackage{stmaryrd}
\usepackage{proof}
\usepackage{booktabs}


\theoremstyle{definition}
\newtheorem{definition}{Définition}[section]
\newtheorem{theorem}{Théorème}[section]
\newtheorem{lemma}[theorem]{Lemme}
\newtheorem*{remark}{Remarque}

\newcommand{\etat}[1]{\langle #1 \rangle}
\newcommand{\etatz}[1]{\langle #1 \rangle_0}
\newcommand{\etatu}[1]{\langle #1 \rangle_1}
\newcommand{\etatd}[1]{\langle #1 \rangle_2}
\newcommand{\etatt}[1]{\langle #1 \rangle_3}
\newcommand{\mcd}{\mathcal{D}}
\newcommand{\mcc}{\mathcal{C}}
\newcommand{\mcl}{\mathcal{L}}
\newcommand{\mtd}{\mathcal{D}}
\newcommand{\mtc}{\mathcal{C}}
\newcommand{\mtl}{\mathcal{L}}
\newcommand{\da}{\downarrow}
\newcommand{\ua}{\uparrow}
\newcommand{\ra}{\rightarrow}
\newcommand{\rsa}{\rightsquigarrow}


\title{Étude d'une mesure d'espace pour le lambda calcul}
\author{Ewen DUFOUR\\[1ex]
{\small Encadré par}\\ {\small Thibaut BALABONSKI}}
\date{Stage de M1 MPRI\\mars--juillet 2025}
\begin{document}

\maketitle

% \begin{abstract}
%   Nous sommes beaux Quelque chose de bien orthographiée.
% \end{abstract}
%
% \tableofcontents
% \section{Introduction}
\section{Le Linear Substitution Calculus}

\begin{definition}[Linear Substitution Calculus]
Les termes du LSC sont donnés par la grammaire~:
\begin{align*}
    t ::= x~|~t~t~|~\lambda x.t~|~t[x / t]
\end{align*}
Les règles de réduction sont~:
\begin{center}
\begin{tabular}{l l l l }
  $\mcl[\lambda x.t]u$ &$\mapsto_{d\beta}$ &$\mcl[t[x / u]]$ & \\
$\mathcal{C}\llbracket x \rrbracket [x / u]$ &$\mapsto_{ls}$ &$\mathcal{C}\llbracket u \rrbracket [x / u]$ & \\
$t[x / u]$ &$\mapsto_{gc}$ &$t$ & if $x\notin fv(t)$
\end{tabular}
\end{center}
où $\mcc$ et $\mcl$ sont des contextes construits par les grammaires~:
\[\mcl ::= \square~|~\mcl[x/t]\]
\[\mcc ::= \square~|~\mcc[x/t]~|~t[x/\mcc]~|~\mcc~t~|~t~\mcc~|~\lambda x.\mcc\]  
\end{definition}
\begin{definition}[Free variables]
\begin{align*}
    fv(x) &= \{ x \} \\
    fv(\lambda x.t) &= fv(t) \setminus \{ x \} \\
    fv(t~u) &= fv(t) \cup fv(u) \\
    fv(t[x/u]) &= (fv(t) \setminus \{ x \}) \cup fv(u) \\
\end{align*}
\end{definition}

\section{LS(G)C}
Nous considérons les termes suivant comme des valeurs~:
\begin{align*}
    v ::= (\lambda x.t)~|~v[x/t]
\end{align*}
\begin{definition}[Contextes d'évaluation]
\begin{align*}
    \mathcal{C} ::= \square~|~\mathcal{C}~t~|~v~\mathcal{C}~|~\mathcal{C}[x/t]
\end{align*}
\end{definition}
\begin{definition}[Contextes de substitution]
\begin{align*}
  \mathcal{L} ::= \square~|~\mathcal{L}[x/v]
\end{align*}
\end{definition}

\begin{definition}[Règles de réduction de LS(G)C nommé]
      \begin{align*}
    &\infer[LA]{t~u\rightarrow t'~u}{t \rightarrow t'}
 &\infer[RA]{v~t\rightarrow v~t'}{t \rightarrow t'}\\
&\infer[GC]{t[x/u]\rightarrow t'}{t \rightarrow t' & x \notin fv(t') }
&\infer[US]{t[x/u]\rightarrow t'[x/u]}{t \rightarrow t' & x \in fv(t')}\\
&\infer[LS]{\mathcal{C}\llbracket x \rrbracket [x/u]\rightarrow \mathcal{C}\llbracket u \rrbracket [x/u]}{x \in fv(\mathcal{C}\llbracket u \rrbracket) }
&\infer[LSGC]{\mathcal{C}\llbracket x \rrbracket [x/u]\rightarrow \mathcal{C}\llbracket u \rrbracket}{x \notin fv(\mathcal{C}\llbracket u \rrbracket) }\\
&\infer[DB]{\mathcal{L}[\lambda x.t]~v \rightarrow \mathcal{L}[t[x/v]]}{x\in fv(t)}
&\infer[DBGC]{\mathcal{L}[\lambda x.t]~v \rightarrow \mathcal{L}[t]}{x\notin fv(t)}
\end{align*}
\end{definition}

\begin{definition}[Règles de réduction de LS(G)C indicé]
      \begin{align*}
    &\infer[LA]{t~u\rightarrow t'~u}{t \rightarrow t'}
 &\infer[RA]{v~t\rightarrow v~t'}{t \rightarrow t'}\\
&\infer[GC]{t[u]\rightarrow \da t'}{t \rightarrow t' & 0 \notin fv(t') }
&\infer[US]{t[u]\rightarrow t'[u]}{t \rightarrow t' & 0 \in fv(t')}\\
&\infer[LS]{\mathcal{C}\llbracket n \rrbracket [u]\rightarrow \mathcal{C}\llbracket \ua^{n+1} u \rrbracket [u]}{0 \in fv(\mathcal{C}\llbracket \ua^{n+1} u \rrbracket)}
&\infer[LSGC]{\mathcal{C}\llbracket n \rrbracket [u]\rightarrow \da \mathcal{C}\llbracket \ua^{n+1} u \rrbracket}{0 \notin fv(\mathcal{C}\llbracket \ua^{n+1} u \rrbracket) }\\
&\infer[DB]{\mathcal{L}[\lambda t]~v \rightarrow \mathcal{L}[t[\ua^{k} v]]}{0 \in fv(t) & k= len(dec(\mcl))}
&\infer[DBGC]{\mathcal{L}[\lambda t]~v \rightarrow \mathcal{L}[\da t]}{0 \notin fv(t)}
\end{align*}
où $\da t$ représente le décrément de $1$ des variables libres de $t$, $\ua$ leur incrément de $1$, 
$dec$ est la fonction qui prend un contexte et le décompose en parties atomiques et $len$ est la fonction qui calcule la longueure de la décomposition.
\end{definition}



\begin{definition}[Taille d'un terme]
  \begin{align*}
    |x| &= |dB(x)|\\
    |\lambda x.t| &= |t| + 1\\
    |t~u| = |t[x/u]| &= |t| + |u|
  \end{align*}
\end{definition}
\begin{remark}
  Nous étendons trivialement cette définitions au différents contextes.
\end{remark}

\begin{definition}[Propreté]
\begin{align*}
    &\infer{cl(x)}{}
 &\infer{cl(\lambda x.t)}{cl(t)}\\
 &\infer{cl(t~u)}{cl(t) & cl(u) }
&\infer{cl(t[x/u])}{cl(t)&cl(u)&x\in fv(t)}
  \end{align*}
Un terme $t$ est dit propre si et seulement si $cl(t)$ possède une dérivation.
\end{definition}

\subsection{Propreté de LS(G)C}
A des fins de clareté nous utilisons ici la version nommée des règles de réduction. Nous pensons que le travail à suivre peut être effectué sur la version indicée d'une manière analogue.
\begin{lemma}[Préservation des variables libres]
  \label{lemmapresfv}
  Soient $\mathcal{C}$ un contexte de d'évaluation, $x$ et $y$ des variables et t un terme.
  \[
    \text{Si}~y \neq x~\wedge~y\in fv(\mathcal{C}[x])~\text{alors}~y \in fv(\mathcal{C}[t])
  \]
\end{lemma}
\begin{proof}
Par induction sur la forme de $\mathcal{C}$
 \begin{itemize}
    \item \underbar{Cas $\mathcal{C} = \square$~:} Prémisse fausse.
    \item \underbar{Cas $\mathcal{C} = \mathcal{C}_1[z/v]$~:}\\
      Comme $y\in fv(\mathcal{C}[x])$, nous avons deux cas.
      \begin{itemize}
        \item \underbar{Cas $y \in fv(v)$~:}\\ 
          En particulier, $y \in fv(\mathcal{C}_1[t][z/v])~\text{c'est à dire}~y \in fv(\mathcal{C}[t])$
        \item \underbar{Cas $y\in fv(\mathcal{C}_1[x]) \wedge y \neq z$~:}\\
          Par hypothèse d'induction, $y\in fv(\mathcal{C}_1[t])$.\\
          Comme $y \neq z$, $y\in fv(\mathcal{C}_1[t][z/v])~\text{i.e}~y \in fv(\mathcal{C}[t]$ 
      \end{itemize}
    \item \underbar{Cas $\mathcal{C} = \mathcal{C}_1~u$~:}\\
      De nouveau, nous avons deux cas.
      \begin{itemize}
        \item \underbar{Cas $y \in fv(u)$~:}\\
          $ y \in fv(\mathcal{C}[t]~u) \Leftrightarrow y \in fv(\mathcal{C}[t])$
        \item \underbar{Cas $y \in fv(\mathcal{C}_1[x])$~:}
          Par hypothèse d'induction, $y \in fv(\mathcal{C}_1[t])$\\ 
          Donc, $y \in fv(\mathcal{C}_1[t]~u)~\text{c.à.d}~y \in fv(\mathcal{C}[t])$
      \end{itemize}
    \item \underbar{Cas $\mathcal{C} = v~\mathcal{C}_1$~:} Analogue au cas précédent.
  \end{itemize} 
\end{proof}

\begin{lemma}[Propreté des substitutions]
  \label{prls}
 Soient $\mathcal{C}$ un contexte d'évaluation, $t$ un terme et $x$ une variable.
 \[ \text{Si}~\mathcal{C}\llbracket x \rrbracket~\text{est propre et}~t~\text{est propre, alors}~ \mathcal{C}[t]~\text{est propre}\]
\end{lemma}
\begin{proof}
  Par induction sur la forme de $\mathcal{C}$.
  \begin{itemize}
    \item \underbar{Cas $\mathcal{C} = \square$~:} trivial.
    \item \underbar{Cas $\mathcal{C} = \mathcal{C}_1[y/v]$~:}\\
      Comme $\mathcal{C}\llbracket x\rrbracket$ est propre, $\mathcal{C}_1\llbracket x \rrbracket$ et $v$ sont propres et $y\in fv(\mathcal{C}_1\llbracket x\rrbracket)$.
      En particulier, $x \neq y$, car $x$ est supposé libre dans $\mathcal{C}\llbracket x \rrbracket$.
      Par conséquence du lemme de préservation des variables libres \ref{lemmapresfv}, $y \in fv(\mathcal{C}_1[t])$.\\
      Par hypothèse d'induction, nous avons $\mathcal{C}_1[t]$ propre.\\
      Donc $\mathcal{C}[t]$ est propre.
    \item \underbar{Cas $\mathcal{C} = \mathcal{C}_1~u$~:}
      Par la propreté de $\mathcal{C}\llbracket x \rrbracket$, nous déduisons que $u$ et $\mathcal{C}_1\llbracket x \rrbracket$ sont propres.
      Notre hypothèse d'induction nous donne $\mathcal{C}_1[t]$ propre. 
      Nous déduisons que $\mathcal{C}[t] = \mathcal{C}_1[t]~u$ est propre.
    \item \underbar{Cas $\mathcal{C} = v~\mathcal{C}_1$~:} Analogue au cas précédent.
  \end{itemize}
\end{proof}

\begin{lemma}
  \label{lemma0}
  Soient $\mathcal{L}$ un contexte de substitution, $t$ un terme et $x$ une variable.
  \[ 
    \text{Si}~x \in fv(\mathcal{L}[\lambda y.t])~\text{, alors}~x \in fv(\mathcal{L}[t])
  \]
\end{lemma}
\begin{proof}
  Par induction sur la forme de $\mathcal{L}$.
  \begin{itemize}
    \item \underbar{Cas $\mathcal{L} = \square$~:} Immédiat.
    \item \underbar{Cas $\mathcal{L}=\mathcal{L}_1[z/v]$~:}\\
      Comme $x\in fv(\mathcal{L}[\lambda y.t]) \Leftrightarrow x \in fv((\mathcal{L}_1[\lambda y.t])[z/v])$, nous avons deux cas.
      \begin{itemize}
        \item \underbar{Cas $x\in fv(v)$~:} Trivialement, $x \in fv((\mathcal{L}_1[t])[z/v]) \Leftrightarrow x\in fv(\mathcal{L}[t])$
        \item \underbar{Cas $x\in fv(\mathcal{L}_1[\lambda y.t])\wedge x \neq z$~:}
          Par hypothèse d'induction, $x\in fv(\mathcal{L}_1[t])$.
          Nous déduisons $x \in fv((\mathcal{L}_1[t])[z/v]) \Leftrightarrow x\in fv(\mathcal{L}[t])$.
      \end{itemize}
  \end{itemize}
\end{proof}

\begin{lemma}[Propreté des applications nettoyante]
  \label{prdbgc}
  Soient $\mathcal{L}$ un contexte de substitution et $t$ un terme.
 \[
   \text{Si}~\mathcal{L}[\lambda x.t]~\text{est propre, alors}~\mathcal{L}[t]~\text{est propre}
 \]
\end{lemma}
\begin{proof}
  Par induction sur la forme de $\mathcal{L}$.
  \begin{itemize}
    \item \underbar{Cas $\mathcal{L} = \square$~:} Par définition de la propreté.
    \item \underbar{Cas $\mathcal{L}=\mathcal{L}_1[y/v]$~:}\\
      Comme $\mathcal{L}[\lambda x.t]~\text{propre}~\Leftrightarrow (\mathcal{L}_1[\lambda x.t])[y/v]~\text{propre}$, nous pouvons en déduire trois choses.
      \begin{enumerate}
        \item $\mathcal{L}_1[\lambda x.t]~\text{propre. Nous en déduisons donc que}~\mathcal{L}_1[t]$~est propre par hypothèse d'induction.
        \item $y\in fv(\mathcal{L}_1[\lambda x.t])$ et par conséquence du lemme \ref{lemma0}, $y \in fv(\mathcal{L}_1[t])$
        \item $v$ propre.
      \end{enumerate}
      Donc $(\mathcal{L}_1[t])[y/v] = \mathcal{L}[t]$~est propre.
  \end{itemize}
\end{proof}
\begin{remark}
  Bien que ce lemme fonctionne lorsque $x \in fv(t)$, il n'a vraiment d'utilité que dans le cas où $x \notin fv(t)$.
\end{remark}

\begin{lemma}
  \label{lemma1}
  Soient $\mathcal{L}$ un contexte de substitution, $t$ et $u$ des termes et $x$ une variable.
  \[ 
    \text{Si}~x \in fv(\mathcal{L}[\lambda y.t])~\text{, alors}~x \in fv(\mathcal{L}[t[y/u]])
  \]
\end{lemma}
\begin{proof}
  Par induction sur la forme de $\mathcal{L}$.
  \begin{itemize}
    \item \underbar{Cas $\mathcal{L} = \square$~:} Obligatoirement, $x\in fv(t)$ et $x\neq y$.
      Par conséquent, $x\in fv(\mathcal{L}[t[y/u]])$.
    \item \underbar{Cas $\mathcal{L}=\mathcal{L}_1[z/v]$~:}\\
      Comme $x\in fv(\mathcal{L}[\lambda y.t])$, nécessairement, $x \in fv((\mathcal{L}_1[\lambda y.t])[z/v])$. Nous avons donc deux cas.
      \begin{itemize}
        \item \underbar{Cas $x\in fv(v)$~:} Trivialement, $x \in fv((\mathcal{L}_1[t[y/u]])[z/v])$. Autrement dit, $x\in fv(\mathcal{L}[t[y/u]])$
        \item \underbar{Cas $x\in fv(\mathcal{L}_1[\lambda y.t])\wedge x \neq z$~:}
          Par hypothèse d'induction, $x\in fv(\mathcal{L}_1[t[y/u]])$.
          Nous déduisons $x \in fv((\mathcal{L}_1[t[y/u]])[z/v]) \Leftrightarrow x\in fv(\mathcal{L}[t[y/u]])$.
      \end{itemize}
  \end{itemize}
\end{proof}

\begin{lemma}[Propreté des applications]
  \label{prdb}
  Soient $\mathcal{L}$ un contexte de substitution et $t$ un terme.
 \[
  \text{Si }x\in fv(t)~\wedge~\mathcal{L}[\lambda x.t]~v~\text{est propre, alors}~\mathcal{L}[t[x/v]]~\text{est propre}
 \]
\end{lemma}
\begin{proof}
  Par induction sur la forme de $\mathcal{L}$.
  \begin{itemize}
    \item \underbar{Cas $\mathcal{L} = \square$~:} Par définition de la propreté, on a $t$ et $v$ propres.\\
      Nous déduisons que $t[x/v] = \mathcal{L}[t[x/v]]$~est propre.
    \item \underbar{Cas $\mathcal{L}=\mathcal{L}_1[y/u]$~:}\\
      Comme $\mathcal{L}[\lambda x.t]~\text{est propre, nécessairement,}~(\mathcal{L}_1[\lambda x.t])[y/u]~$est propre. Nous pouvons en déduire trois choses.
      \begin{enumerate}
        \item $\mathcal{L}_1[\lambda x.t]~\text{est propre, et par hypothèse d'induction,}~\mathcal{L}_1[t[x/v]]$~est propre également.
        \item $y\in fv(\mathcal{L}_1[\lambda x.t])~\text{ce qui nous donne}~y \in fv(\mathcal{L}_1[t[x/v]])$~par le lemme \ref{lemma1}.
        \item $u$ propre.
      \end{enumerate}
      Donc $(\mathcal{L}_1[t[x/v]])[y/u] = \mathcal{L}[t[x/v]]$~est propre.
  \end{itemize}
\end{proof}

\begin{lemma}[Préservation de la propreté]
  Soient $t_0$ et $t_1$ des termes. 
  \[
    t_0~\text{propre}~\wedge~t_0\rightarrow~t_1~\Rightarrow~t_1~\text{propre}
  \]
\end{lemma}
\begin{proof}
  Par induction sur la dérivation de $t_0 \rightarrow t_1$.
  \begin{itemize}
    \item \underbar{Cas DB~:} Conséquence directe du lemme \ref{prdb}. 
    \item \underbar{Cas DBGC~:} Conséquence du lemme \ref{prdbgc}. 
    \item \underbar{Cas LS et LSGC~:} Conséquences du lemme de propreté des substitutions \ref{prls}.
    \item \underbar{Cas GC~:} Supposons que $t[x/u]$ est propre. Donc, $t$ est propre. Or, $t \rightarrow t'$. Par hypothèse d'induction, $t'$ est propre.
    \item \underbar{Cas US~:} Supposons que $t[x/u]$ est propre. Il en découle, $t$ et $u$ sont propres. Or, $t \rightarrow t'$. Par hypothèse d'induction, $t'$ est propre.
      Comme dans les hypothèses d'US, $x \in fv(t')$, $t'[x/u]$ est propre.
    \item \underbar{Cas LA~:} Supposons $t~u$ propre. Comme $t \rightarrow t'$, nous déduisons que $t'$ est propre et par conséquence, $t'~u$ est propre.
    \item \underbar{Cas RA~:} Analogue au cas précédent.
  \end{itemize}
\end{proof}

\section{FLAME}
FLAME est une machine abstraite qui implante la stratégie d'appel par valeur pour les termes du LSC. 
Un état de FLAME se décrit par un triplet annoté de la forme $\etat{(v,n),t,\mcc}_e$ où~:
\begin{itemize}
  \item $t$ représente le terme en train d'être évalué.
  \item $\mcc$ est une pile de contexte d'évaluation.
    \item $(v, n)$ contient optionnellement l'argument d'une application et sa \og distance\fg~par rapport à l'abstraction.
    \item $e$ est un nombre qui indique la phase de fonctionnement de FLAME.
    \end{itemize} 
Toutes les transitions de FLAME sont décrites dans la figure~\ref{FLAME} en annexe.

\begin{definition}[Taille d'un état de FLAME]
\begin{align*}
    |\etat{(v,n), t, \mcc}_e| = |t| + |v| + |n| + |\mcc|
\end{align*}  
\end{definition}


\subsection{Simulation du LSC en appel par valeur par FLAME}
Dans la suite de ce travail, nous utilisons la version indicée de LS(G)C.
\begin{lemma}[Décomposition]
  \label{lemmaplouf}
  Soient $\mathcal{C}$, $\mathcal{D}$ des contextes d'évaluation et $t$ un terme.
\[\langle -,\mathcal{C}[t], \mathcal{D}\rangle_0 \rightsquigarrow^* \langle -,t,\mathcal{C} :: \mathcal{D}\rangle_0\]
en utilisant un espace $|\mcc|+|t|+|\mcd|$
\end{lemma}
\begin{proof}
  Par induction sur la forme de $\mathcal{C}$.
  \begin{itemize}
    \item \underbar{Cas $\mathcal{C} = \square$~:} immédiat.\\
    \item \underbar{Cas $\mathcal{C} = \mathcal{C}_1[v]$~:}\\
  \begin{align*}
    \langle -, \mathcal{C}[t], \mathcal{D}\rangle_0 = &\langle -, \mathcal{C}_1[t][v], \mathcal{D}\rangle_0 \\
    \rightsquigarrow_{us} &\langle -, \mathcal{C}_1[t], \square[v]:: \mathcal{D}\rangle_0 \\
    \rightsquigarrow^* &\langle -, t, \mathcal{C}_1:: \square[v]:: \mathcal{D}\rangle_0 \quad \text{Par hypothèse d'induction.}
  \end{align*}
    \item \underbar{Cas $\mathcal{C} = \mathcal{C}_1~u$~:}\\
  \begin{align*}
    \langle -, \mathcal{C}[t], \mathcal{D}\rangle_0 = &\langle -, \mathcal{C}_1[t]~u, \mathcal{D}\rangle_0 \\
    \rightsquigarrow_{lapp} &\langle -, \mathcal{C}_1[t], \square~u:: \mathcal{D}\rangle_0 \\
    \rightsquigarrow^* &\langle -, t, \mathcal{C}_1:: \square~u:: \mathcal{D}\rangle_0 \quad \text{Par hypothèse d'induction.}
  \end{align*}
    \item \underbar{Cas $\mathcal{C} = v~\mathcal{C}_1$~:}\\
  \begin{align*}
    \langle -, \mathcal{C}[t], \mathcal{D}\rangle_0 = &\langle -, v~\mathcal{C}_1[t], \mathcal{D}\rangle_0 \\
    \rightsquigarrow_{lapp} &\langle -, \mathcal{C}_1[t], v~\square:: \mathcal{D}\rangle_0 \\
    \rightsquigarrow^* &\langle -, t, \mathcal{C}_1:: v~\square:: \mathcal{D}\rangle_0 \quad \text{Par hypothèse d'induction.}
  \end{align*}
  \end{itemize}
  L'espace utilisé se voit sur les réductions.
\end{proof}

\begin{lemma}[Reconstruction]
  \label{lemmaportugal}
  Soient $\mathcal{C}$ et $\mathcal{D}$ des contextes d'évaluation et $t$ un terme.
  \[\mathcal{C}[t]~\text{propre}~\Rightarrow \langle-,t,\mathcal{C}::\mathcal{D}\rangle_2 \rightsquigarrow^* \langle-,\mathcal{C}[t],\mathcal{D}\rangle_2\]
en utilisant un espace $|\mcc|+|t|+|\mcd|$
\end{lemma}
\begin{proof}
  Par induction sur la forme de $\mathcal{C}$.
  \begin{itemize}
    \item \underbar{Cas $\mathcal{C} = \square$~:} immédiat.
    \item \underbar{Cas $\mathcal{C} = \mathcal{C}_1[v]$~:}\\
      Nous avons
       $ \mathcal{C}[t]~\text{propre,}~\text{c.à.d}~\mathcal{C}_1[t][v]~\text{propre, ce qui nous donne}~0\in fv(\mathcal{C}_1[t])~\wedge~\mathcal{C}_1[t]~$propre.
      Et par conséquent~:
      \begin{align*}
        \langle -, t, \mathcal{C} :: \mathcal{D}\rangle_2 = &\langle -, t, \mathcal{C}_1[v]:: \mathcal{D}\rangle_2 \\
        = &\langle -, t, \mathcal{C}_1:: \square[v] :: \mathcal{D}\rangle_2 \\
        \rightsquigarrow^* &\langle -, \mathcal{C}_1[t], \square[v]:: \mathcal{D}\rangle_2 \quad &\text{Par hypothèse d'induction.}\\
        \rightsquigarrow_{bsubst} &\langle -, \mathcal{C}_1[t][v], \mathcal{D}\rangle_2 \quad &\text{Car}~0\in fv(\mathcal{C}_1[t])
      \end{align*}
    \item \underbar{Cas $\mathcal{C} = \mathcal{C}_1~u$~:}\\
      Nous avons $\mathcal{C}[t]~\text{propre, c'est à dire}~\mathcal{C}_1[t]~u~\text{propre. Cela force}~\mathcal{C}_1[t]~$propre.
      Ce qui nous donne~:
      \begin{align*}
        \langle -, t, \mathcal{C} :: \mathcal{D}\rangle_2 = &\langle -, t, \mathcal{C}_1~u:: \mathcal{D}\rangle_2 \\
        = &\langle -, t, \mathcal{C}_1:: \square~u :: \mathcal{D}\rangle_2 \\
        \rightsquigarrow^* &\langle -, \mathcal{C}_1[t], \square~u:: \mathcal{D}\rangle_2 \quad &\text{Par hypothèse d'induction.}\\
        \rightsquigarrow_{blapp} &\langle -, \mathcal{C}_1[t]~u, \mathcal{D}\rangle_2 \quad &\text{Par propreté de}~\mathcal{C}
      \end{align*}
    \item \underbar{Cas $\mathcal{C} = v~\mathcal{C}_1$~:} Analogue au cas précédent.
  \end{itemize}
  L'espace utilisé se voit sur les réductions.
\end{proof}


\begin{lemma}[Plongeon substitution]
  \label{lemmaplouf2}
  Soient $\mathcal{D}$ un contexte d'évaluation, $\mcl$ un contexte de substitution et $t$, $u$ des termes.
  \[ \etatu{(u,n), \mcl[t], \mcd} \rightsquigarrow_{rdb}^k \etatu{(u,n+k), t, \mcl::\mcd} \]
  où $k$ est la longueur de la décomposition de $\mcl$,
en utilisant un espace $|\mcd|+|t|+|\mcl|+|u| + |n+k|$
\end{lemma}
\begin{proof}
  Par induction sur la forme de $\mathcal{L}$.
  \begin{itemize}
    \item \underbar{Cas $\mathcal{L} = \square$~:} immédiat.\\
    \item \underbar{Cas $\mathcal{L} = \mathcal{L}_1[v]$~:}\\
  \begin{align*}
    \langle (u,n), \mathcal{L}[t], \mathcal{D}\rangle_1 = &\langle u, \mathcal{L}_1[t][v], \mathcal{D}\rangle_1 \\
    \rightsquigarrow_{rdb} &\langle (u,n+1), \mathcal{L}_1[t], \square[v]:: \mathcal{D}\rangle_1 \\
  \rightsquigarrow^{k-1}_{rdb} &\langle (u,n+k), t, \mathcal{L}_1:: \square[v]:: \mathcal{D}\rangle_1 \quad \text{Par hypothèse d'induction.}\\
 = &\etatu{(u,n+k), t, \mcl_1[v]::\mcd} = \etatu{(u,n+k), t, \mcl::\mcd}
  \end{align*}
  \end{itemize}
  L'espace utilisé se voit sur les réductions.
\end{proof}

\begin{lemma}[Simulation raisonnable]
  \label{lem:sim}
  Soient $t$ et $u$ des termes propres et $\mcd$ un contexte d'évaluation.
  \[
    \text{Si}~t \rightarrow u~\text{avec un espace}~S~\text{, alors }~\etat{-,t,\mcd}_0 \rightsquigarrow^* \etat{-,u,\mcd}_2~\text{avec un espace maximum}~2 \times S+|\mcd|
  \]
\end{lemma}
\begin{proof}
  Par induction sur la dérivation de $t\rightarrow u$.
  \begin{itemize}
    \item \underbar{Cas LS~:} Nous avons $\mtc\llbracket n\rrbracket[v]\rightarrow \mcc[\ua^{n+1} v][v]$ sous hypothèse que $0\in fv(\mcc[\ua^{n+1} v])$. Donc $S = |\mcc|+|v| + max(|n|, |\ua^{n+1} v|)$.\\
      \begin{tabular}{@{}l c c c@{}}
        \toprule
        transition & état & taille & justification\\
        \midrule
                   & $\etatz{-,\mcc\llbracket n\rrbracket[v],\mcd}$ & $|\mcd|+|v|+|\mcc|+|n|$\\
        $\rightsquigarrow_{us}$ & $ \etatz{-, \mcc\llbracket n \rrbracket, \square[v]::\mcd}$ & $|\mcd| + |v| + |\mcc|+|n|$\\ 
        $\rightsquigarrow^*$ & $ \etatz{-, n ,\mcc:: \square[v]::\mcd}$ & $|\mcd| + |v| + |\mcc|+|n|$& Lemme \ref{lemmaplouf}\\
        $\rightsquigarrow_{ls}$ & $ \etatd{-, \ua^{n+1} v ,\mcc:: \square[v]::\mcd}$ & $|\mcd| + |\ua^{n+1} v| + |\mcc|+|v|$\\
        $\rightsquigarrow^*$ & $ \etatd{-, \mcc[\ua^{n+1} v], \square[v]::\mcd}$ & $|\mcd| + |\ua^{n+1} v| + |\mcc|+|v|$ &Lemme \ref{lemmaportugal}\\ 
        $\rightsquigarrow_{bsubst}$ & $ \etatd{-, \mcc[\ua^{n+1} v][v], \mcd}$ & $|\mcd| + |\ua^{n+1}v| + |\mcc|+|v|$&$0\in fv(\mcc[\ua^{n+1} v])$\\ 
        \bottomrule
      \end{tabular}\\
      La taille maximum utilisée est $S +|\mcd|$.
    \item \underbar{Cas LSGC~:} Nous avons $\mtc\llbracket n\rrbracket[v]\rightarrow \da\mcc[\ua^{n+1} v]$ avec comme hypothèse que $0\notin fv(\mcc[\ua^{n+1} v])$. Donc $S = |\mcc|+|v| + max(|n|, |\ua^{n+1} v|)$.\\
      \begin{tabular}{@{}l c c c@{}}
        \toprule
        transition & état & taille & justification\\
        \midrule
                   & $\etatz{-,\mcc\llbracket n\rrbracket[v],\mcd}$ & $|\mcd|+|v|+|\mcc|+|n|$\\
        $\rightsquigarrow_{us}$ & $ \etatz{-, \mcc\llbracket n \rrbracket, \square[v]::\mcd}$ & $|\mcd| + |v| + |\mcc|+|n|$\\ 
        $\rightsquigarrow^*$ & $ \etatz{-, n ,\mcc:: \square[v]::\mcd}$ & $|\mcd| + |v| + |\mcc|+|n|$& Lemme \ref{lemmaplouf}\\
        $\rightsquigarrow_{ls}$ & $ \etatd{-, \ua^{n+1} v ,\mcc:: \square[v]::\mcd}$ & $|\mcd| + |v| + |\mcc|+|\ua^{n+1} v|$\\
        $\rightsquigarrow^*$ & $ \etatd{-, \mcc[\ua^{n+1} v], \square[v]::\mcd}$ & $|\mcd| + |v| + |\mcc| +|\ua^{n+1} v|$ &Lemme \ref{lemmaportugal}\\ 
        $\rightsquigarrow_{gc}$ & $ \etatd{-, \da \mcc[\ua^{n+1} v], \mcd}$ & $|\mcd| + |\da \mcc[\ua^{n+1} v]|$& $0\notin fv(\mcc[\ua^{n+1} v])$\\ 
        \bottomrule
      \end{tabular}\\
      La taille maximum utilisée est $S+|\mcd|$.
    \item \underbar{Cas DB~:} Nous avons $\mcl[\lambda t]~v \rightarrow \mcl[t[\ua^{k} v]]$ sous hypothèse que $0\in fv(t)$. Donc $S = |\mcl|+|t|+max(1+|v|, |\ua^k v|)$.\\
      \begin{tabular}{@{}l c c c@{}}
        \toprule
        transition & état & taille & justification\\
        \midrule
                   & $\etatz{-,\mcl[\lambda t]~v,\mcd}$ & $|\mcd|+|v|+|\mcl|+|t|+1$\\
        $\rightsquigarrow_{db}$ & $ \etatu{(v,0), \mcl[ \lambda t], \mcd}$ & $|\mcd| + |v| + |\mcl|+|t|+1$\\ 
        $\rightsquigarrow^k_{rdb}$ & $ \etatu{(v,k), \lambda t ,\mtl::\mtd}$ & $|\mtd| + |v| + |\mtl|+|t|+|k|+1$& Lemme \ref{lemmaplouf2}\\
        $\rightsquigarrow_{fdb}$ & $ \etatd{-, t ,\square[\ua^k v]::\mtl::\mtd}$ & $|\mtd| + |\ua^k v| + |\mtl|+|t|$\\
        $\rightsquigarrow_{rsubst}$ & $ \etatd{-, t[\ua^k v] ,\mtl::\mtd}$ & $|\mtd| +|\ua^k v|+ |\mtl|+|t|$& $0 \in fv(t)$\\
        $\rightsquigarrow^*$ & $ \etatd{-, \mtl[t[\ua^k v]],\mtd}$ & $|\mtd| + |v| + |\mtl|+|t|$ &Lemme \ref{lemmaportugal}\\ 
        \bottomrule
      \end{tabular}\\
      Le lemme \ref{lemmaplouf2} nous assure que la transition de la ligne deux vers la ligne trois ce fait dans un espace~:
     \begin{align*} 
       |\mtd|+ |v| + |\mtl|+|t|+|k|+1 &\le |\mtd| + |v| + |\mtl|+|t|+|\mtl|+1 \\
        &\le 2 \times (|\mtl|+|t|+|v|+1)+|\mcd|\\ 
        &= 2 \times S + |\mcd|
     \end{align*}
      De plus, les autres transitions se font en espace $S+|D|$.
    \item \underbar{Cas DBGC~:} Nous avons $\mtl[\lambda t]~v \rightarrow \mtl[t]$ sous hypothèse que $0\notin fv(t)$. Donc $S = |\mcl|+|t|+max(1+|v|, |\ua^k v|)$.\\
      \begin{tabular}{@{}l c c c@{}}
        \toprule
        transition & état & taille & justification\\
        \midrule
                  & $\etatz{-,\mcl[\lambda t]~v,\mcd}$ & $|\mcd|+|v|+|\mcl|+|t|+1$\\
        $\rightsquigarrow_{db}$ & $ \etatu{(v,0), \mcl[ \lambda t], \mcd}$ & $|\mcd| + |v| + |\mcl|+|t|+1$\\ 
        $\rightsquigarrow^k_{rdb}$ & $ \etatu{(v,k), \lambda t ,\mtl::\mtd}$ & $|\mtd| + |v| + |\mtl|+|t|+|k|+1$& Lemme \ref{lemmaplouf2}\\
        $\rightsquigarrow_{fdb}$ & $ \etatd{-, t ,\square[\ua^k v]::\mtl::\mtd}$ & $|\mtd| + |\ua^k v| + |\mtl|+|t|$\\
        $\rightsquigarrow_{gc}$ & $ \etatd{-, t ,\mtl::\mtd}$ & $|\mtd| + |\mtl|+|t|$& $x\notin fv(t)$\\
        $\rightsquigarrow^*$ & $ \etatd{-, \mtl[t],\mtd}$ & $|\mtd| + |\mtl|+|t|$ &Lemme \ref{lemmaportugal}\\ 
        \bottomrule
      \end{tabular}\\
      Une analyse similaire au celle du cas précédent nous assure que la réduction se fait dans un espace maximum $2\times S +|\mtd|$.
    \item \underbar{Cas GC~:} Nous avons $t[x/v] \rightarrow \da t'$ sous hypothèse que $t \rightarrow t'$ et $0\notin fv(t')$. Donc $S = S'+|v|$. Nous notons $S' = |t \rightarrow t'| $.\\
      \begin{tabular}{@{}l c c c@{}}
        \toprule
        transition & état & taille & justification\\
        \midrule
                   & $\etatz{-,t[v],\mtd}$ & $|\mtd|+|v|+|t|$\\
        $\rightsquigarrow_{us}$ & $ \etatz{-, t, \square[v]::\mtd}$ & $|\mtd| + |v| + |t|$\\ 
        $\rightsquigarrow^*$ & $ \etatd{-, t' ,\square[v]::\mtd}$ & $|\mtd| + |v| + |t'|$& Par hypothèse d'induction\\
      $\rightsquigarrow_{gc}$ & $ \etatd{-, \da t' ,\mtd}$ & $|\mtd| + |\da t'|$&$ 0\notin fv(t') $\\
        \bottomrule
      \end{tabular}\\
      Notre hypothèse d'induction indique également que la transition de la deuxième ligne vers la troisième ligne se fait avec un espace: 
      \begin{align*}
        2 \times S' +|\mtd| + |v| &\le 2\times(S'+|v|) +|\mtd|\\ 
                                  &= 2 \times S + |\mtd|
       \end{align*}
Comme $|t| \le S'$ et $|\da t'|\le|t'| \le S'$, les autres transitions s'effectuent dans un espace $S+|\mtd|$.
      % Nous devons distinguer trois cas.
      % \begin{itemize}
      %   \item \underbar{Cas $S = S' = |t'|$~:}
      %     \begin{align*}
      %       k \times S' + |\mtd| + |v|  &\le k \times S'+|\mtd|+|t'|  \quad \text{par définition de $S$}\\
      %       &= (k+1) \times S' +|\mtd| \\ 
      %       &= (k+1)\times S +|\mtd|
      %     \end{align*}
      %   \item \underbar{Cas $S' = |t|$ et $S = |t| + |v|$~:}
      %     \begin{align*}
      %       k \times S' + |v| + |\mtd| &\le k \times (S' + |v|) + |\mtd| \\
      %                                  & = k \times S + |\mtd|
      %     \end{align*}
      %   \item \underbar{Cas $S' = |t'|$ et $S = |t| + |v|$~:}
      %     \begin{align*}
      %       k \times S' + |\mtd| + |v|  &=  k \times |t'|+|\mtd|+|v|\\
      %                                   &= k \times S + |v| +|\mtd| \quad \text{Car $S \ge |t'|$ par définition.} \\ 
      %       &= (k+1)\times S +|\mtd|
      %     \end{align*}
      % \end{itemize}
    \item \underbar{Cas US~:} Nous avons $t[v] \rightarrow t'[v]$ sous hypothèse que $0\in fv(t')$ et  $t\rightarrow t'$. Donc $S = S'+|v|$ avec $S' = | t \rightarrow t'|$.\\
      \begin{tabular}{@{}l c c c@{}}
        \toprule
        transition & état & taille & justification\\
        \midrule
                   & $\etatz{-,t[x/v],\mcd}$ & $|\mcd|+|v|+|t|$\\
        $\rightsquigarrow_{us}$ & $ \etatz{-, t, \square[x/v]:: \mcd}$ & $|\mcd| + |v| +|t|$\\ 
        $\rightsquigarrow^{*}$ & $ \etatd{-, t' ,\square[x/v]::\mcd}$ & $|v| + |\mcd|+|t'|$& Par hypothèse d'induction \\
        $\rightsquigarrow^{*}$ & $ \etatd{-, t'[x/v], \mcd}$ & $|v| + |\mcd|+|t'|$&$x \in fv(t')$\\
        \bottomrule
      \end{tabular}\\
      Une analyse similaire au celle du cas précédent nous assure que la réduction se fait dans un espace maximum $2\times S +|\mtd|$.
    \item \underbar{Cas LA~:} Nous avons $t~u \rightarrow t'~u$ sous hypothèse que $t\rightarrow t'$. Donc $S = S'+|u|$ avec $S' = | t \rightarrow t'|$.\\
      \begin{tabular}{@{}l c c c@{}}
        \toprule
        transition & état & taille & justification\\
        \midrule
                   & $\etatz{-,t~u,\mtd}$ & $|\mtd|+|u|+|t|$\\
        $\rightsquigarrow_{lapp}$ & $ \etatz{-, t, \square~u:: \mtd}$ & $|\mtd| + |u| +|t|$\\ 
        $\rightsquigarrow^*$ & $ \etatz{-, t' ,\square~u::\mtd}$ & $|u| + |\mtd|+|t'|$& Par hypothèse d'induction\\
        $\rightsquigarrow^*$ & $ \etatz{-, t'~u, \mtd}$ & $|u| + |\mtd|+|t'|$\\
        \bottomrule
      \end{tabular}\\
      Notre hypothèse nous affirme que la transition de la deuxième ligne vers la troisième ligne se fait avec un espace $ 2 \times S' +|\mtd| + |u|$ . 
      La taille maximum utilisée est $2 \times S' +|\mtd|+|u| \le 2 \times (S'+|u|)+|\mtd| = 2 \times S + |\mtd|$.
    \item \underbar{Cas RA~:} Analogue au cas précédent.
  \end{itemize}
\end{proof}

\begin{theorem}[FLAME simule LS(G)C]
  Soient $t$ un terme propre et $\rho$ sa séquence de réduction. 
  \begin{enumerate}
    \item Si $\rho$ diverge, alors la réduction de $\etatz{-, t, \square}$ diverge. 
    \item Si $\rho$ termine sur une valeur $u$, alors $\etatz{-, t, \square} \rsa^{*} \etatt{-, u, \square}$
    \item L'espace pris par la réduction du point 2 est $2 \times |\rho| = 2 \times |t \ra^* u|$
  \end{enumerate}

\end{theorem}
\begin{proof}
  Par application successive du lemme \ref{lem:sim}.
\end{proof}

\subsection{Simulation de FLAME par les machines de Turing}

\section{Le Space Calculus}

\begin{definition}[Termes du Space Calculus]
Le Space calcul est défini par les termes donnés par la grammaire suivante~:
\begin{align*}
    l ::&= x \\
    &|~\lambda x.a \\
    &|~a~a\\
    a ::&= l^{\alpha}\\
    c ::&= a\\
        &|~c[x/c]\\
    t ::&= c\\
        &|~t~c
\end{align*}
ainsi que par les règles de réductions~:
\begin{align*}
&\infer[sea_v]{((t^\beta~x^\gamma)^\alpha e) \rightarrow (t^\beta e_{|t})~e(x)}{}
&\infer[sea_{\neg v}]{((t^\beta~u^\gamma)^\alpha e) \rightarrow (t^\beta e_{|t})~(u^\gamma e_{|u}) }{}\\\\
&\infer[\beta_{\neg w}]{((\lambda x .t^\beta)^\alpha e_t)~(u^\gamma e_u ) \rightarrow (t^\beta[x/u^\gamma e_u]e_t)}{x \in fv(t) }
&\infer[\beta_w]{((\lambda x .t^\beta)^\alpha e_t)~(u^\gamma e_u ) \rightarrow (t^\beta e_t)}{x \notin fv(t) }\\\\
&\infer[sub]{x^\alpha[x/u^\beta e] \rightarrow u^\beta e}{}
&\infer[cont]{\mathcal{C}[t] \rightarrow \mathcal{C}[t']}{t\rightarrow t'}
\end{align*}
où $e_{|t}$ dénote la restriction de l'environnement $e$ aux variables libres de $t$ et 
$\mathcal{C}$ est un contexte de réduction de la forme~:
\begin{align*}
  \mathcal{C} ::= \square~|~\mathcal{C}~c.
\end{align*}

\end{definition}
\begin{definition}[Restriction d'un environnement]
\begin{align*}
a_{|t} &= a\\
    (e[x/c])_{|t} &= e_{|t}[x/c] \quad &\text{si } x\in fv(t)\\
    (e[x/c])_{|t} &= e_{|t} \quad &\text{si } x\notin fv(t)
\end{align*}
\end{definition}

\section*{Annexe}
\begin{figure}[h]
\label{FLAME}
\centering
\resizebox{\textwidth}{!}{%
\begin{tabular}{@{}l l l l l l l l l l@{}}
\toprule
\textbf{arg} & \textbf{terme} & \textbf{contexte} & \textbf{état} & \textbf{transition}& \textbf{arg} & \textbf{terme} & \textbf{contexte} & \textbf{état} \\
\midrule

– & $t[ u ] $ & $\mcc$ & 0 & $\rsa_{us}$ & – & $t$ & $\textsf{LS}(u) :: \mcc$ & 0  \\

- & $t~u$ & $\mcc$ & 0 & $\rsa_{lapp}$ & – & $t$ & $\textsf{RA}(u) ::\mcc$ & 0  \\
- & $v~t$ & $\mcc$ & 0 & $\rsa_{rapp}$ & – & $t$ & $\textsf{LA}(v) :: \mcc$ & 0  \\
- & $n$  & $\mcc$ & 0 & $\rsa_{ls}$ & - & $\uparrow^{n+1} t$ & $\mcc$ & 2 & $\textsf{LS}(t)=\textsf{nth\_subst}(\mcc)$ \\
- & $v_1~v_2$  & $\mcc$ & 0 & $\rsa_{db}$ & $(v_2,0)$ & $v_1$ & $\mcc$ & 1  \\

$(v_2, n)$ & $v_1[ t ]$ &$\mcc$ & 1 & $\rsa_{rdb}$ & $(v_2, n+1)$ & $v_1$ & $\textsf{LS}(t) :: \mcc$ & 1\\ 
$(v, n)$ & $\lambda.t $ &$\mcc$ & 1 & $\rsa_{fdb}$ & - & $t$ & $\textsf{LS}(\uparrow^{n}v) :: \mcc$ & 2 & $0\in fv(t)$\\
-&$v$ & $[]$ & 2 & $\rsa_{stop}$ & -& $v$ &$[]$ & 3 \\
-&$t$ & $[]$ & 2 & $\rsa_{dive}$ & -& $t$ &$[]$ & 0 \\
-&$t$ & $\textsf{LS}(u) :: \mcc$ & 2 & $\rsa_{gc}$ & -& $\downarrow t$ &$\mcc$ & 2 & $0\notin \textsf{fv}(t)$ \\
-&$t$ & $\textsf{LS}(u) :: \mcc$ & 2 & $\rsa_{bsubst}$ & -& $t[ u ]$ &$\mcc$ & 2 & $0\in \textsf{fv}(t)$ \\
-&$t$ & $\textsf{RA}(u) :: \mcc$ & 2 & $\rsa_{blapp}$ & -& $t~u$ &$\mcc$ & 2 \\
-&$t$ & $\textsf{LA}(u) :: \mcc$ & 2 & $\rsa_{brapp}$ & -& $u~t$ &$\mcc$ & 2 \\

\bottomrule
\end{tabular}
}
\caption{Transitions de FLAME.}
\end{figure}
\end{document}

